% ===== PRÉAMBULE CANONIQUE — à coller VERBATIM en tête de CHAQUE .tex =====
\documentclass[11pt,a4paper]{article}
\usepackage[utf8]{inputenc}\usepackage[T1]{fontenc}\usepackage{lmodern}
\usepackage[french]{babel}
\usepackage{amsmath,amssymb,mathtools}\usepackage{icomma}
\usepackage[version=4]{mhchem}          % \ce{...} : équations & formules
\usepackage{siunitx}\sisetup{locale=FR} % \qty{}{}, \si{}, \ang{}
\usepackage{chemfig}                    % \chemfig{...} : structures développées
\usepackage{xcolor}\usepackage{colortbl}
\usepackage{tikz}\usepackage{pgfplots}\pgfplotsset{compat=1.18}
\usepackage[most]{tcolorbox}\usepackage{enumitem}\usepackage{array,booktabs}
\usepackage[a4paper,margin=2.2cm]{geometry}
\usetikzlibrary{arrows.meta,calc,angles,quotes,positioning,patterns,backgrounds,shapes.geometric}
\definecolor{primary}{RGB}{222,49,99}\definecolor{primarydark}{RGB}{150,22,55}
\definecolor{defcol}{RGB}{0,114,178}\definecolor{methcol}{RGB}{52,92,148}
\definecolor{excol}{RGB}{0,135,90}\definecolor{attncol}{RGB}{200,40,55}
\tcbset{boxbase/.style={enhanced,breakable,boxrule=0.9pt,arc=2.5pt,
  left=3mm,right=3mm,top=2.3mm,bottom=2.3mm,fonttitle=\bfseries,coltitle=white}}
% --- boîtes TUTORIEL (noms EXACTS, ne pas renommer) ---
\newtcolorbox{cle}[1][]{boxbase,colback=defcol!6,colframe=defcol,
  title={\textsc{À retenir}\ifx\relax#1\relax\relax\else\textmd{ : \itshape #1}\fi}}
\newtcolorbox{methode}[1][]{boxbase,colback=methcol!7,colframe=methcol,sharp corners,
  title={\textsc{Méthode}\ifx\relax#1\relax\relax\else\textmd{ --- \itshape #1}\fi}}
\newtcolorbox{exemplebox}[1][]{boxbase,colback=excol!5,colframe=excol,
  title={\textsc{Exemple}\ifx\relax#1\relax\relax\else\textmd{ : \itshape #1}\fi}}
\newtcolorbox{piege}[1][]{boxbase,colback=attncol!6,colframe=attncol,
  title={\textsc{Piège}\ifx\relax#1\relax\relax\else\textmd{ : \itshape #1}\fi}}
\newtcolorbox{remarque}[1][]{boxbase,colback=black!4,colframe=black!45,
  title={\textsc{Remarque}\ifx\relax#1\relax\relax\else\textmd{ : \itshape #1}\fi}}
% --- boîtes EXERCICE / CORRIGÉ (noms EXACTS, auto-comptées) ---
\newtcolorbox[auto counter]{exo}[1][]{boxbase,colback=primary!4,colframe=primary,
  title={\textsc{Exercice}~\thetcbcounter\ifx\relax#1\relax\relax\else\textmd{ --- \itshape #1}\fi}}
\newtcolorbox[auto counter]{corrige}[1][]{boxbase,colback=excol!4,colframe=excol,
  title={\textsc{Corrigé}~\thetcbcounter\ifx\relax#1\relax\relax\else\textmd{ --- \itshape #1}\fi}}
\newcommand{\R}{\mathbb{R}}\newcommand{\N}{\mathbb{N}}\newcommand{\Z}{\mathbb{Z}}\newcommand{\ds}{\displaystyle}
\begin{document}
\begin{corrige}[Classer tous les carbones d'une chaîne ramifiée]
Numérotons la chaîne principale $\ce{C1H3-C2H(-C5H3)-C3H2-C4H3}$ (le $\ce{C5}$ est le méthyle ramifié).
\begin{itemize}[leftmargin=1.6em,itemsep=1pt]
  \item $\ce{C1}$ ($\ce{CH3}$) : $1$ voisin carboné $\Rightarrow$ \textbf{primaire}.
  \item $\ce{C2}$ : lié à $\ce{C1}$, $\ce{C3}$ et $\ce{C5}$, soit $3$ voisins $\Rightarrow$ \textbf{tertiaire}.
  \item $\ce{C3}$ ($\ce{CH2}$) : $2$ voisins $\Rightarrow$ \textbf{secondaire}.
  \item $\ce{C4}$ ($\ce{CH3}$) : $1$ voisin $\Rightarrow$ \textbf{primaire}.
  \item $\ce{C5}$ (méthyle) : $1$ voisin $\Rightarrow$ \textbf{primaire}.
\end{itemize}
\textbf{Décompte :} $3$ primaires, $1$ secondaire, $1$ tertiaire, $0$ quaternaire.
\end{corrige}

\begin{corrige}[Dénombrer les isomères de constitution]
\begin{enumerate}[label=\alph*)]
  \item $\ce{C4H10}$ : \textbf{2} isomères --- le \textbf{butane} (chaîne droite) et le \textbf{2-méthylpropane} (isobutane).
  \item $\ce{C5H12}$ : \textbf{3} isomères --- le \textbf{pentane}, le \textbf{2-méthylbutane} (isopentane) et le \textbf{2,2-diméthylpropane} (néopentane).
\end{enumerate}
\end{corrige}

\begin{corrige}[Qui présente l'isomérie cis/trans ?]
\begin{enumerate}[label=\alph*)]
  \item \textbf{Oui.} Chaque carbone de la double liaison porte un $\ce{CH3}$ et un $\ce{H}$ (deux groupes différents). Les deux configurations sont :
  \begin{center}
  \begin{tabular}{cc}
  \chemfig{CH_3-[:30](-[:90]H)=[:30](-[:90]H)-[:-30]CH_3} &
  \chemfig{CH_3-[:30](-[:90]H)=[:-30](-[:90]H)-[:30]CH_3}\\
  \emph{cis} (les deux $\ce{CH3}$ du même côté) & \emph{trans} (de part et d'autre)
  \end{tabular}
  \end{center}
  \item \textbf{Non.} Sur le but-1-ène, le carbone terminal $\ce{CH2=}$ porte \emph{deux} atomes d'hydrogène (deux groupes \emph{identiques}). La condition « deux groupes différents sur chaque carbone de la $\ce{C=C}$ » n'est pas remplie : pas d'isomérie cis/trans.
\end{enumerate}
\end{corrige}

\begin{corrige}[Lire une formule topologique]
\begin{enumerate}[label=\alph*)]
  \item Le squelette a $4$ sommets sur la chaîne principale $+$ $1$ branche $=$ \textbf{5 carbones}.
  \item Alcane à $5$~C $\Rightarrow \ce{C5H12}$.
  \item Une ramification méthyle sur le deuxième carbone : c'est le \textbf{2-méthylbutane}.
\end{enumerate}
\end{corrige}

\begin{corrige}[Compter les H de chaque carbone (molécule insaturée)]
\ce{CH2=CH-CHO}, de gauche à droite :
\begin{itemize}[leftmargin=1.6em,itemsep=1pt]
  \item $\ce{CH2=}$ : une double liaison C (compte $2$) $\Rightarrow n_{\text{H}}=4-2=\textbf{2}$~H.
  \item $=\ce{CH}-$ : une double ($2$) $+$ une simple ($1$) $=3$ $\Rightarrow n_{\text{H}}=4-3=\textbf{1}$~H.
  \item $-\ce{CHO}$ : une simple vers C ($1$) $+$ une double vers O ($2$) $=3$ $\Rightarrow n_{\text{H}}=4-3=\textbf{1}$~H.
\end{itemize}
Total : $2+1+1=4$~H, cohérent avec la formule brute $\ce{C3H4O}$.
\end{corrige}

\end{document}
